% Bench measurement and comparison with the post-layout simulation. Macros from meas.tex
% (written by make_report.py from the run*.uart.log of ../ro_run.sh and pex_summary.csv of
% ro_reliability/sim/pex/run_all.sh).
\section{Measurement}
\label{sec:meas}
Logs \MeasRuns{} of \texttt{ro\_run.sh} (firmware \texttt{ro\_test}, \texttt{ro\_diag}, \texttt{ro\_scan}),
RO supplies from the on-board \SI{1.2}{V} regulator through JP3/JP4, room temperature. The versions of
the firmware, the repository and the bitstream of every run are in \texttt{versions.txt}.

\begin{center}\small
\begin{tabular}{lcc}
\toprule
 & RO13 & RO101 \\
\midrule
Pad, ring alone (\SI{1}{s} gate) & \SI{\MeasAloneA}{MHz} & \SI{\MeasAloneB}{MHz} \\
Pad, both rings on & \SI{\MeasBothA}{MHz} & \SI{\MeasBothB}{MHz} \\
Ring (pad $\times$ 32), alone & \SI{\MeasRingAloneA}{MHz} & \SI{\MeasRingAloneB}{MHz} \\
Change when the other ring runs & \SI{\MeasBothShiftA}{\percent} & \SI{\MeasBothShiftB}{\percent} \\
Repeatability (both on, std.\ dev.) & \SI{\MeasBothSdA}{\percent} & \SI{\MeasBothSdB}{\percent} \\
Spread over \MeasRestartsA{} restarts (\SI{100}{ms} gates) & \SI{\MeasRestartSpreadA}{\percent} & \SI{\MeasRestartSpreadB}{\percent} \\
Duty cycle of the pad & \SI{\MeasDutyA}{\percent} & \SI{\MeasDutyB}{\percent} \\
\bottomrule
\end{tabular}
\end{center}
The ratio $f_{13}/f_{101}$ is \MeasRatio, against \RatioAB{} predicted.

\subsection{Function}
\begin{itemize}
  \item The SPI writes arrive: the rings start and stop with the configuration (Table~\ref{tab:scan}).
        They oscillate only with \texttt{RO\_control} = 1, \texttt{n\_RO\_control} = 0 and
        \texttt{DUT\_Footer} = 1, as designed.
  \item \texttt{RO\_RST\_B} low stops both outputs, and the outputs have a \SI{50}{\percent} duty cycle:
        the counts come from the $\div 32$ dividers.
  \item \texttt{DUT\_Header} (bit 3) has no visible function: with the header off (\texttt{0x0e}) the
        rings keep running (RO13 \SI{\HdrOffShiftA}{\percent}, RO101 \SI{\HdrOffShiftB}{\percent}),
        where the DUT can only pull down and the ring should stop. In the extracted cells the header
        gate is on the \texttt{DUT\_Header} pin, so the suspect is the chip-level connection from
        \texttt{R[3]}/\texttt{R[11]}.
  \item With \texttt{RO\_control} = \texttt{n\_RO\_control} = 0 and \texttt{DUT\_gate} = 1
        (\texttt{0x04}, \texttt{0x0c}) RO13 shows a slow oscillation of about \SI{10}{kHz} at the pad;
        RO101 does not. Not understood; outside the oscillation mode.
  \item \texttt{DUT\_gate}, disconnected from the ring in oscillation mode, still moves RO13 by about
        \SI{0.6}{\percent} (coupling through the off transmission gate) and RO101 by \SI{0.05}{\percent}.
\end{itemize}

\begin{table}[H]\centering\scriptsize
\input{table_scan.tex}
\caption{Configuration scan: pad edges per second of both rings for every value of the 4 bits of
the RO register (written to both rings), from a ring-open state.}
\label{tab:scan}
\end{table}

\ifdefined\HavePex
\section{Comparison with the post-layout simulation}
\label{sec:cmp}
The two cells were cut out of the fabricated chip GDS and extracted with KLayout-PEX 2.5D, coupling
capacitances included (\texttt{ro\_reliability/sim/pex}); the extraction keeps the 13- and 101-stage
loops (checked on the extracted netlists). They are simulated in the conditions of the bench, with the
pad of \texttt{extra\_load}, from the restart of the firmware (the footer off, then on, which launches a
single edge).
\begin{table}[H]\centering\small
\input{table_cmp.tex}
\caption{Ring frequency: simulation against the board (pad $\times$ 32, ring alone).}
\label{tab:cmpcases}
\end{table}

\subsection{RO13}
The wiring of the layout lowers RO13 by \SI{\PexWireA}{\percent}, to \SI{\PexA}{MHz} at TT; the board
gives \MeasOverPexA{} of that, below even the SS corner (\SI{\PexSsA}{MHz}). The post-layout ring would
run at the measured frequency with a supply of about \SI{\EquivVdd}{V} (Figure~\ref{fig:vddpex}).
The extraction has no resistance and the transistors are typical, so part of the difference may be
the chip, but the bench points at the supply as well: RO13 slows by \SI{\MeasBothShiftA}{\percent}
when RO101 also runs, which at about \SI{12}{\percent} per \SI{100}{mV} means some \SI{40}{mV} lost in
the supply path for the \SI{0.5}{mA} of the other ring. Both RO supplies come from the same regulator
through JP3 and JP4, and both rings share one ground pad.
\begin{figure}[H]\centering
\includegraphics{fig/vdd_pex.pdf}
\caption{Post-layout RO13 against its supply, and the measured frequency.}
\label{fig:vddpex}
\end{figure}

\subsection{RO101}
RO101 cannot run at \SI{\MeasRingAloneB}{MHz} with one edge in the ring: that is \SI{37}{ps} per stage,
and the post-layout fundamental is \SI{\PexB}{MHz}. Scaled by the factor RO13 shows
(\MeasOverPexA), its fundamental would be \SI{\RoBFundEst}{MHz}, and the measurement is
\RoBEdges{} times that: the ring would hold about seven edges at once (a harmonic mode), each still
divided by 32 at the pad.

The simulation does not keep such a mode. Started from all nodes at zero, the extracted ring does
begin with several edges (\SI{\StartPexFirst}{MHz} of edges on \texttt{extra\_load}, \StartPexEdges{}
edges), but they catch up with each other and vanish within \SI{\StartPexSettle}{\micro s}, and the
ring ends at its fundamental (\SI{\PexUicB}{MHz}). Without the wiring capacitance the same start gives
\StartDevEdges{} edges, gone within \SI{\StartDevSettle}{\micro s}. The board, on the other hand, always
lands on the same frequency, over \MeasRestartsB{} restarts made in six different ways, including the
footer-off restart that launches a single edge in simulation. So either the chip keeps creating edges
until about seven circulate, or the ring that oscillates is not the 101-stage loop: seven edges in 101
stages run at the same rate as one edge in a loop of about 14 stages, close to RO13's 13. The extracted
cell has a single 101-stage loop (\texttt{ringloop.py}), so a short loop would have to come from outside
the cell or from a defect. The \texttt{extra\_load} node, slow because of its pad, is the first suspect
for extra edges: a slow edge with noise on it can switch the next stage twice.
\fi

\section{Reproducing this report}
\label{sec:repro}
Everything in this report comes from two repositories, both from scripts; the numbers in the text are
macros computed from the data files by \texttt{make\_report.py}.
\begin{description}
  \item[Bench] (RISCVConsole) SoC with the \texttt{rocounter} device (\texttt{fpga/ULX3S}:
        \texttt{make verilog bitstream program}, tools from \texttt{\textasciitilde/init\_openlane.sh}),
        firmware \texttt{software/RISCVConsoleCode} (\texttt{make elf}, \texttt{src/ro.c}), wiring in
        \texttt{software/ro/README.md}. \texttt{software/ro/ro\_run.sh <dir>} boots the firmware over
        JTAG and stores \texttt{runN.uart.log} and the versions of firmware, repository and bitstream
        in \texttt{versions.txt}.
  \item[Schematic] (\texttt{asicone\_sg13g2/ro\_reliability}) \texttt{python3 sim/ro\_sim.py board}
        $\rightarrow$ \texttt{sim.csv}; ngspice-47 with OSDI \cite{ngspice}, IHP-Open-PDK \cite{ihppdk}.
        \texttt{sim/spi/run.sh} checks the SPI frame of the firmware against the chip RTL.
  \item[Post-layout] (same repository) \texttt{sim/pex/run\_all.sh}: cuts \texttt{top\_13} and
        \texttt{top\_101} out of the chip GDS, LVS and KLayout-PEX 2.5D, checks the ring loops,
        simulates the cases of Table~\ref{tab:cmpcases} $\rightarrow$ \texttt{pex\_summary.csv}, and
        the start-up of the all-zero starts $\rightarrow$ \texttt{pex\_startup.csv}
        (about \SI{5}{h} on 32 cores).
  \item[Report] copy the three CSV files into \texttt{software/ro/results/<dir>} next to the logs, then
        \texttt{make -C software/ro/report RESULTS=../results/<dir>}.
\end{description}
This report: \texttt{results/\ResultsDir}.

\section{Next steps}
\begin{enumerate}
  \item Measure ROVDD and RO2VDD (JP3, JP4) with a multimeter, with each ring off, alone and both on.
  \item Feed ROVDD from a bench supply and sweep it: in a harmonic mode the frequency of RO101 jumps by
        discrete ratios (7, 5, 3, 1 edges), while RO13 changes smoothly; a short loop would change smoothly
        too, following RO13.
  \item Look at \texttt{ro\_101\_ex} (J1 pin 28) with a scope, short ground: several pulses per period
        instead of one square wave would confirm the harmonic mode.
  \item Check the chip-level route of \texttt{R[3]}/\texttt{R[11]} to \texttt{DUT\_Header} (extraction of
        the chip around the RO macros).
\end{enumerate}
